\documentclass[11pt]{article}

\usepackage[T1]{fontenc}
\usepackage[utf8]{inputenc}
\usepackage{lmodern}
\usepackage{amsmath,amssymb}
\usepackage[letterpaper,margin=1in]{geometry}
\usepackage{enumitem}
\usepackage{hyperref}

\hypersetup{colorlinks=true,linkcolor=blue,urlcolor=blue}
\setlength{\emergencystretch}{3em}
\setlist[itemize]{leftmargin=1.5em,itemsep=2pt,topsep=3pt}
\setlist[enumerate]{leftmargin=1.8em,itemsep=4pt,topsep=3pt}

\title{Preguntas de investigación para el proyecto de carteras\\
de mínima varianza en acciones mexicanas}
\author{Julio César Galindo López\\
\small Universidad Panamericana, Ciudad UP
\and
Laura Jiménez Casillas\\
\small Universidad Panamericana, Ciudad UP}
\date{Borrador de discusión para Laura\\1 de octubre de 2026}

\begin{document}
\maketitle

\begin{abstract}
\noindent
Este documento propone varias preguntas de investigación para discutir y
priorizar entre los coautores antes de fijar el diseño de un artículo
empírico sobre carteras de mínima varianza en acciones mexicanas. Las
alternativas se distinguen por su estimando y por los datos adicionales que
requieren; se recomienda elegir una pregunta principal y tratar las demás
como extensiones sólo si los datos lo permiten.
El punto de partida es un panel Bloomberg de retornos totales brutos para 37
acciones, con 3,706 observaciones diarias entre 2011 y 2025. En los resultados
preliminares disponibles, RIE reduce levemente la volatilidad frente a la
covarianza muestral en la partición fija, pero su Sharpe es menor; la cartera
long-only $1/N$ obtiene el mayor Sharpe puntual tanto en esa partición como en
el ejercicio rolling. Estos resultados son descriptivos de un universo
retrospectivo y una trayectoria histórica. El diseño no resuelve el sesgo de
supervivencia ni incorpora costos de transacción observados. Por ello, el
documento presenta alternativas de pregunta, una propuesta de estructura y
un programa de trabajo; no es una versión lista para envío.
\end{abstract}

\section{Propósito de esta nota}

El objetivo es que los coautores acuerden primero la pregunta, el alcance y
las condiciones mínimas para seguir desarrollando el proyecto. La versión
existente contiene una reestimación inicial y un artículo más desarrollado;
esta nota reduce la propuesta a un plan de revisión. No presupone que Laura
haya aprobado el enfoque, el calendario, la autoría final ni la pregunta
principal.

\section{Preguntas de investigación candidatas}

\subsection{Opción de partida}

\begin{quote}
\emph{¿Cuánto cambia la evaluación fuera de muestra de los estimadores de
covarianza para carteras de mínima varianza en acciones mexicanas cuando se
imponen restricciones long-only, se incluye el benchmark $1/N$ y se considera
rebalanceo periódico?}
\end{quote}

Esta opción aprovecha directamente los análisis ya realizados. Su foco es la
comparación entre estimadores bajo distintas reglas de inversión; su
limitación principal es que los resultados actuales dependen de un universo
retrospectivo y de una sola trayectoria histórica.

\subsection{Alternativas para discutir}

\begin{enumerate}
\item \textbf{¿La limpieza espectral mejora la medición del riesgo que enfrenta
un inversionista?} Comparar, fuera de muestra, las predicciones de matrices de
covarianza de SM, CL, LW y RIE con medidas de riesgo realizadas posteriores.
El estimando central sería la precisión de pronóstico de covarianzas o de
volatilidad de cartera, no el Sharpe de una cartera optimizada. Requiere
definir una función de pérdida antes de observar los resultados y conservar
ventanas rolling sin solapamientos indebidos.

\item \textbf{¿Las restricciones long-only cambian la utilidad práctica de
los estimadores?} Medir cómo varían desempeño, riesgo, concentración y
turnover al pasar de pesos sin restricciones a long-only, manteniendo el
benchmark $1/N$. Esta pregunta separa el efecto del estimador del efecto de
las reglas de inversión. Es la extensión más próxima al análisis existente.

\item \textbf{¿Las carteras optimizadas conservan una ventaja después de
costos de negociación y límites de liquidez?} Comparar rendimientos netos
con costos calculados a partir de spreads y volumen, incluyendo restricciones
de participación y activos difíciles de negociar. Esta pregunta requiere
datos observados de ejecución o proxies locales defendibles; los escenarios
hipotéticos actuales no bastan para responderla.

\item \textbf{¿Cuándo, si acaso, cambia el desempeño relativo según las
condiciones del mercado?} Evaluar si los resultados varían entre estados
definidos ex ante, como periodos de alta y baja volatilidad o liquidez. Los
umbrales y las series que definen los estados deben especificarse sin usar
los resultados de cartera para escogerlos; las pruebas requieren potencia y
observaciones suficientes por estado.

\item \textbf{¿Se generalizan los resultados a otros mercados emergentes?}
Replicar un protocolo armonizado en otros países y comparar patrones de
estimación, restricciones, costos y benchmarks. Es una extensión sustantiva,
pero exige universos punto-en-tiempo, retornos totales y calendarios
comparables; sin esa armonización, las diferencias entre países serían
difíciles de interpretar.
\end{enumerate}

Las alternativas implican objetos distintos y no conviene incorporarlas
todas en una sola pregunta. Para la conversación inicial, se recomienda
elegir una como pregunta principal, fijar su resultado primario y decidir qué
otras preguntas son extensiones condicionadas a la disponibilidad de datos.

\section{Estado de la evidencia}

El panel de trabajo usa el campo Bloomberg
\texttt{TOT\_RETURN\_INDEX\_GROSS\_DVDS} para 37 acciones mapeadas desde una
lista candidata anterior de 38 valores. FUNO11 no tuvo una pestaña utilizable
en el archivo y no se sustituyó. Las series balanceadas abarcan del 18 de
julio de 2011 al 29 de septiembre de 2025. El corte cronológico fijo reserva
741 observaciones para prueba (28 de noviembre de 2022--29 de septiembre de
2025); la ventana previa contiene 2,965 observaciones.

En el holdout fijo con pesos sin restricciones, la volatilidad diaria de RIE
es 0.6952\%, frente a 0.6998\% para SM. El Sharpe anualizado es $-0.232$ para
RIE y $-0.207$ para SM; la diferencia RIE--SM es $-0.0248$, con IC bootstrap
de 95\% $[-0.0451,-0.0038]$.

Con restricciones long-only en el mismo holdout, el Sharpe puntual es
$-0.035$ para RIE y 0.450 para $1/N$. La diferencia RIE--$1/N$ es $-0.485$;
su IC bootstrap de 95\% $[-1.068,0.053]$ incluye cero.

En el análisis rolling, que usa 25 bloques de 21 observaciones y una ventana
estimadora de 504 observaciones, el Sharpe long-only es $-0.646$ para RIE y
0.767 para $1/N$. La diferencia RIE--$1/N$ es $-1.413$, con IC bootstrap de
95\% $[-2.333,-0.641]$. Este análisis comparte la trayectoria histórica con
el holdout fijo y no constituye una réplica independiente.

Finalmente, bajo un escenario hipotético de 50 puntos base por unidad de
turnover unidireccional, el retorno medio anualizado rolling es 0.57\% para
RIE y 18.03\% para $1/N$. No son costos observados; se excluyen costos
iniciales y otras fricciones.

Estos resultados no demuestran que RIE sea generalmente inferior ni que
$1/N$ vaya a superar a las carteras optimizadas fuera de esta muestra.
Describen estimaciones condicionadas a una selección de activos, una
trayectoria, una tasa libre de riesgo fijada en 8\% y pesos estimados. Además,
la significancia nominal del contraste primario no incorpora incertidumbre
por selección del universo o por elección de modelos.

\section{Datos adicionales necesarios}

La prioridad es corregir primero la selección del universo; obtener más
observaciones para las mismas acciones no elimina por sí solo el posible
sesgo de supervivencia. La lista siguiente separa los datos necesarios para
fortalecer el estudio mexicano de los que permitirían sostener una
contribución regional más amplia.

La prioridad principal es reconstruir el \textbf{universo histórico
punto-en-tiempo de BMV/BIVA}. Necesitamos los constituyentes y las fechas
efectivas de entrada y salida, incluidas empresas deslistadas, fusiones y
cambios de ticker, además de identificadores estables —por ejemplo, ISIN u
otro identificador histórico enlazable—. Esto permitiría definir qué activos
podían pertenecer al universo en cada fecha y reducir el sesgo de
supervivencia y selección retrospectiva. Una lista de emisoras actuales no
sería suficiente.

En segundo lugar, hacen falta \textbf{retornos totales históricos para el
universo reconstruido}, incluidas las empresas deslistadas y con tratamiento
consistente de dividendos, splits, derechos y rendimientos de salida. La
frecuencia y cobertura temporal deben ser coherentes; también habría que
documentar observaciones faltantes, acciones corporativas y salidas de
mercado. Estos datos son los que permitirían volver a estimar las carteras
con el conjunto elegible en cada momento.

Para evaluar implementación, sería muy valioso conseguir \textbf{medidas
históricas de liquidez y ejecución}: bid--ask spreads, volumen negociado,
días sin operación y turnover y, si es viable, costos o impacto de mercado.
Deben tener fechas e identificadores compatibles con los retornos. Así se
podrían reemplazar o calibrar los escenarios hipotéticos actuales y evaluar
si las carteras son implementables.

Para interpretar el desempeño se necesita además una \textbf{serie mexicana
fechada de tasa libre de riesgo}, con vencimiento definido —por ejemplo,
CETES adecuado al horizonte—, y un \textbf{benchmark mexicano de retorno
total} cuya metodología y composición estén documentadas. Esto evitaría
depender sólo de la tasa fija de 8\% y situaría las carteras frente al mercado
local.

Por último, para evaluar si los resultados se extienden más allá de México,
convendría buscar datos comparables de uno o más mercados emergentes:
universos históricos, retornos totales, activos deslistados, benchmarks,
tasas libres de riesgo y medidas de liquidez, con reglas armonizadas y fechas
compatibles. Una comparación regional con coberturas distintas podría añadir
sesgos en vez de resolverlos.

\subsection{Orden práctico de adquisición}

\begin{enumerate}
\item Confirmar si es posible obtener membresía histórica BMV/BIVA con
deslistadas e identificadores estables, y bajo qué licencia.
\item Verificar que existan retornos totales utilizables para todos esos
activos y para todo su periodo de elegibilidad. Si no se puede reconstruir
este panel, mantener explícitamente el estudio como evidencia condicionada a
la selección observada.
\item Identificar disponibilidad de spreads, volumen y medidas de liquidez
antes de formular afirmaciones sobre costos o implementación.
\item Añadir una tasa libre de riesgo y benchmark mexicanos como controles de
interpretación; evaluar la extensión a otros países sólo si se pueden
armonizar universo, retornos y calendarios.
\end{enumerate}

\noindent\textbf{Metadatos que deben solicitarse junto con los datos:}
diccionario de campos y unidades, moneda, calendario, reglas de dividendos y
acciones corporativas, fechas efectivas de membresía, tratamiento de
deslistadas y política de uso, publicación y redistribución. Los datos
Bloomberg u otros datos licenciados no se deben incorporar al repositorio
público; antes de compartir resultados agregados también hay que comprobar
que su publicación esté permitida por la licencia.

\section{Esqueleto del artículo}

\begin{enumerate}
\item \textbf{Introducción.} Plantear el problema económico, la pregunta
única y la diferencia entre reducir varianza y mejorar desempeño ajustado por
riesgo. Delimitar desde el inicio que la evidencia es mexicana y no una
estimación de un efecto común a los mercados emergentes.
\item \textbf{Literatura y predicciones.} Revisar la optimización de
media-varianza, error de estimación, limpieza espectral, shrinkage, carteras
long-only y el benchmark $1/N$. Formular predicciones contrastables sin
presuponer que la limpieza domina.
\item \textbf{Datos y universo.} Describir proveedor, campo de retorno total,
fechas, ticker mapping, tratamiento de acciones corporativas y faltantes.
Explicar la exclusión de FUNO11, la selección retrospectiva y el riesgo de
supervivencia. Precisar qué puede y qué no puede redistribuirse por licencia.
\item \textbf{Diseño empírico.} Definir estimadores, optimización sin
restricciones y long-only, benchmark $1/N$, métricas, tasa libre de riesgo,
costos, turnover, ventanas y reglas de inferencia. Separar el test fijo del
walk-forward y evitar presentarlos como muestras independientes.
\item \textbf{Resultados.} Presentar primero la comparación primaria
preespecificada; después, restricciones, benchmark, rolling, turnover y
costos. Reportar estimaciones e intervalos, no sólo rankings puntuales.
\item \textbf{Robustez y límites de diseño.} Evaluar ventanas, bloques de
bootstrap, tasa libre de riesgo, costos plausibles, concentración y universo.
Indicar con claridad qué robusteces son sensibilidad sobre la misma historia.
\item \textbf{Discusión.} Explicar qué implicaciones son defendibles para una
decisión de inversión en México y qué evidencia adicional sería necesaria
para generalizar a otros mercados.
\item \textbf{Conclusión.} Responder la pregunta dentro del alcance de los
datos; enumerar como trabajo futuro un universo punto-en-tiempo, costos
observados y una comparación regional.
\end{enumerate}

\section{Programa de trabajo propuesto}

El calendario es una propuesta de ocho semanas desde el acuerdo de los
coautores. Las fechas y la distribución de responsabilidades deben confirmarse
con Laura antes de tratarse como compromisos.

Durante la \textbf{primera semana}, la conversación con Laura debería cerrar la
pregunta, el estimando primario, la audiencia, el papel de cada autor y los
límites de los datos. El resultado será una lista compartida de decisiones y
preguntas secundarias.

En las \textbf{semanas segunda y tercera}, se auditarán procedencia y universo:
tickers, dividendos, fechas, empresas excluidas y posibilidad de recuperar
membresía histórica con deslistadas. El producto será un diccionario de datos
y un mapa de cobertura, junto con una decisión explícita sobre ampliar el
universo o mantener el análisis como evidencia exploratoria condicionada.

Entre las \textbf{semanas tercera y cuarta}, se preespecificará la validación:
walk-forward, reglas de reestimación, comparadores, ventanas, métricas e
inferencia primaria. También se verificará el cálculo de turnover y costos.
La salida será un protocolo reproducible y comprobaciones de consistencia
entre código, tablas y texto.

En las \textbf{semanas cuarta y quinta}, se buscará la factibilidad de obtener
costos de ejecución, medidas de liquidez y un benchmark local; además, se
probará la sensibilidad a la tasa libre de riesgo y a los escenarios de
costos. Si no hay datos de ejecución viables, los escenarios se conservarán
como hipotéticos y se etiquetarán con claridad.

Durante las \textbf{semanas quinta y sexta}, se completarán la estimación ampliada,
la inferencia, los gráficos y las sensibilidades de estabilidad temporal.
El objetivo es congelar un paquete de resultados sin modificar la pregunta
primaria a la luz de los resultados.

En las \textbf{semanas sexta y séptima}, se redactará el manuscrito en
coautoría y se revisará que la pregunta, la contribución y el alcance estén
alineados con la evidencia ampliada. La salida será un borrador completo con
limitaciones y audiencia académica acordadas.

En la \textbf{octava semana}, se realizará la auditoría final y la réplica desde
los datos licenciados, se revisarán referencias y disponibilidad de datos, y
se decidirá si conviene cerrar el manuscrito, ampliar el proyecto o cambiar
su alcance.

\subsection{Decisiones prioritarias para conversar con Laura}

\begin{enumerate}
\item ¿La pregunta y la contribución propuesta responden a un problema que
ambos autores consideran publicable e interesante?
\item ¿Se puede obtener un universo histórico punto-en-tiempo, incluyendo
deslistadas, y bajo qué licencia y esfuerzo?
\item ¿Hay acceso viable a información local de spreads, volumen, liquidez o
costos de ejecución para sustituir o validar los escenarios hipotéticos?
\item ¿Qué proveedor puede entregar el universo BMV/BIVA punto-en-tiempo, con
deslistadas e identificadores estables, y los retornos totales asociados?
\item ¿Hay datos comparables de otros mercados emergentes con definiciones y
cobertura temporal compatibles, si se decide explorar una comparación
entre países?
\item ¿Qué resultados deben ser confirmatorios y cuáles quedarían claramente
como exploratorios?
\end{enumerate}

\section{Criterios para continuar y riesgos}

\begin{itemize}
\item \textbf{Fijar una pregunta principal} sólo después de acordar cuál
estimando responde a un problema empírico claro y qué datos permiten
identificarlo.
\item \textbf{Ampliar el diseño} si se puede reconstruir el universo histórico
y obtener datos de costos o liquidez; de lo contrario, presentar los
resultados como evidencia condicionada y limitar las afirmaciones.
\item \textbf{No presentar} los costos hipotéticos como costos mexicanos
observados, el rolling como réplica independiente, ni el universo balanceado
como libre de sesgo de supervivencia.
\item \textbf{Proteger datos licenciados}: compartir sólo resultados
agregados permitidos y código que no redistribuya series, precios ni pesos
por activo.
\end{itemize}

\section{Siguiente paso inmediato}

Usar esta nota para una reunión de alcance con Laura. Si ambos acuerdan
cuál de las preguntas candidatas priorizar y qué datos son factibles, el
primer entregable será el protocolo empírico preespecificado. Si no hay
acuerdo, revisar las alternativas antes de invertir en una nueva corrida o
en una reescritura extensa.

\vfill
\noindent\textbf{Documento de discusión.} Los resultados incluidos son
preliminares y corresponden al paquete agregado de trabajo disponible al
1 de octubre de 2026. Este proyecto no se presenta como artículo aceptado ni
como manuscrito listo para envío.

\end{document}
